\chapter{Extension 3: the real ocean and the sea-level constant}
\label{ch:ocean}

\section{The problem}

The core model puts water everywhere; the plate of Chapter~\ref{ch:plate} puts it only in the oceans for the restoring force.
But the water that follows the geoid also has to be confined to the oceans: on land, a rise of the geoid does not add water.
Once that is so, the water's load and mass depend on $C(x)\,N(x)$, which couples degrees, and $N$ itself depends on the water.
So the geoid and the water must be solved together.

A second problem appears with real land and sea: what fixes absolute sea level?
In the core model only anomalies matter.
With a reference lithosphere and continents, the mean level of the ocean relative to the columns determines which areas are land and how high they stand, and a wrong mean produces a spurious land--ocean step.

\section{How we tackle the ocean}

The sea surface is an equipotential, so it follows the geoid: where the masses of the model raise the geoid by $N$ at sea, the water is $N$ deeper; on land nothing happens.
This is the static version of the sea-level equation of glacial isostatic adjustment, in which the ocean also follows the geoid it helps to create.

To first order the ocean that follows the geoid does two things at sea and nothing on land:
\begin{itemize}[nosep]
\item it loads the plate with its weight, $\delta q = -\rho_w g\,C N$;
\item it adds its mass to the geoid as a sheet at sea level, $\delta N = S(R)\,\rho_w C N$.
\end{itemize}
Writing $r$ for the geoid the model computes before the ocean is added (the direct geoid of all masses, plus the geoid of the deflection under their load), the total geoid satisfies
\begin{equation}
    N = r + S(R)\,\rho_w C N - G_w P\,(\rho_w g\,C N)
    \quad\Longleftrightarrow\quad
    (I + T)\,N = r,
    \qquad T = \rho_w\,(g\,G_w P - S(R))\,C,
\end{equation}
with $P$ the plate (load to deflection) and $G_w$ the geoid of the deflected column.
Both effects depend on $N$ and both change it, which is why the equation is implicit.
The water's own load and its own mass nearly cancel, as in the core model: the plate sinks under the water's weight by almost enough to remove an equal mass from the surface. So $T$ is small: its spectral radius is $0.0014$ at $L=50$.
We solve by a Neumann series, $N = \sum_k (-T)^k r$, which converges in about four terms.
The deflection is then $w = P(q - \rho_w g C N)$ and the observed topography, measured from sea level on land and at sea, $h = w - N$.

\begin{defence}
\Q{Is the self-gravitating ocean worth the trouble?}
\A The water sheet is about \qty{7}{\percent} of the mantle geoid at degree 2, though its weight sinks the plate and cancels most of that in $N$; its larger effects are on $w$ (about \qty{2}{\percent} at degree 2) and through $h=w-N$, with $N/w$ about \qty{5}{\percent} at degree 2.
Confining the water to the oceans also matters, because on land there is none.
An earlier version that took the water term from flowkernels' free-surface geoid, applied it on land as well, and never added the water's mass to the geoid had the mantle geoid \qty{11}{\percent} low at degree 2.
The fix is cheap and the limit test confirms it to $10^{-10}$.
\end{defence}

\section{How we tackle sea level}

Every column has an absolute mass per unit area, $M(x)$, from its surface to the lid base, ice included.
Heavy columns sit low and light columns stand high, but relative to what?

\begin{figure}[h]
\centering
\begin{tikzpicture}[x=1cm,y=-1cm,font=\small]
\fill[gray!15] (0,3.0) rectangle (10,5.2);
\fill[blue!15] (3.6,1.2) rectangle (5.6,1.8);
\fill[teal!25] (0.6,1.2) rectangle (2.6,4.2);
\fill[teal!25] (3.6,1.8) rectangle (5.6,4.8);
\fill[teal!25] (6.6,0.6) rectangle (8.6,3.6);
\draw (0.6,1.2) rectangle (2.6,4.2); \draw (3.6,1.8) rectangle (5.6,4.8); \draw (6.6,0.6) rectangle (8.6,3.6);
\draw[blue!70!black, dashed] (0,1.2) -- (10,1.2) node[right, font=\scriptsize] {sea level};
\draw[dashed] (0,5.0) -- (10,5.0) node[right, font=\scriptsize] {depth $Z$};
\node at (1.6,2.6) {$M = K$}; \node[font=\scriptsize] at (1.6,3.0) {at sea level};
\node at (4.6,3.2) {$M > K$}; \node[font=\scriptsize] at (4.6,3.6) {heavy: ocean};
\node at (7.6,2.0) {$M < K$}; \node[font=\scriptsize] at (7.6,2.4) {light: land};
\node[font=\scriptsize] at (4.6,1.5) {water, $-w$ thick};
\draw[-{Stealth}] (5.9,1.25) -- (5.9,1.75) node[midway,right,font=\scriptsize] {$w<0$};
\draw[-{Stealth}] (8.9,1.15) -- (8.9,0.65) node[midway,right,font=\scriptsize] {$w>0$};
\node[font=\scriptsize] at (5,5.45) {asthenosphere $\rho_{\mathrm{LAB}}$; equal pressure at $Z$; all columns of length $T$};
\end{tikzpicture}
\caption{The column balance behind the sea-level constant.}
\end{figure}

Compare columns of equal length on the same asthenosphere, and require equal pressure at a depth $Z$ below them (force up = force down).
Under a column of mass $M$ whose surface is at height $w$ above sea level (positive up), the mass above $Z$ is
\[
    \underbrace{M}_{\text{column}} + \underbrace{\rho_w C\,(-w)}_{\text{water above}} + \underbrace{\rho_{\mathrm{LAB}}\,(Z - T + w)}_{\text{asthenosphere below}} .
\]
At sea $w<0$ and the water fills from the seafloor at $w$ up to sea level at 0, so its thickness is $-w$; on land $C=0$ removes the term.
Equating with the reference column ($M = K$, $w = 0$) and cancelling $\rho_{\mathrm{LAB}}(Z-T)$:
\begin{equation}
    (\rho_{\mathrm{LAB}}-\rho_w C)\,w = K - M .
\end{equation}
$K$ is a single global number, the mass a column must have to float with its surface exactly at sea level.
The per-column quantity is $K - M(x)$, the column's mass deficit or excess, and $w = (K-M)/(\rho_{\mathrm{LAB}}-\rho_w C)$ converts it into a height.

A change $\Delta K$ is not a uniform shift: it moves columns by $\Delta K/2340$ at sea and $\Delta K/3370$ on land, because water fills an oceanic deflection but not a continental one.
So a wrong $K$ creates a step between land and ocean.
In the core model, with water everywhere, the contrast is uniform and the constant $b_0$ absorbs any change of $K$, so it is not needed there.
The reference load is $q_{\mathrm{ref}} = -g(M - \langle M\rangle) + gK_0$: the mean column mass plus a constant.

\begin{decision}{Fit the sea-level constant on columns that hold their whole lithosphere}
\textbf{Choice.} $K_0$ is fitted, by least squares on the fine grid, to the observed bedrock elevation, using only columns whose WINTERC-G LAB is shallower than the column base.
A perturbation $\delta K$, with prior standard deviation \qty{1e6}{kg/m^2}, is an unknown of the inversion.

\textbf{Alternatives.} $K = \langle M\rangle$ (zero-mean load); fitting on all columns; conserving ocean mass.

\textbf{Why.}
\begin{itemize}[nosep]
\item With $K = \langle M\rangle$ there was a \qty{705}{m} spurious land--ocean step and \qty{316}{m} rms leaking into all degrees.
\item Where the LAB is deeper than the column base, the dense cold root is missing from $M$, so those columns are biased; fitting on them would push the constant to compensate.
\item After the fit, thin-lithosphere oceans match the seafloor to \qty{2}{m}, and continents stand too high in proportion to their LAB depth (\qty{+0.07}{km} for LAB $<\qty{100}{km}$, rising to \qty{+0.87}{km} beyond \qty{200}{km}): exactly the missing roots, left for the mantle unknown.
\end{itemize}
The fitted value is $K_0 = \qty{-5.2e6}{kg/m^2}$ relative to $\langle M\rangle$, equivalent to the mean column standing about \qty{2.2}{km} below sea level with the ocean contrast ($K_0/(\rho_{\mathrm{LAB}}-\rho_w)$).

\textbf{Under (b).} With the lid base following the LAB, cratons get their keels back and $K_0$ can be fitted on every column, which removes the restriction above.

\textbf{Caveat.} $\delta K$ is nearly degenerate with the degree-0 Moho displacement (correlation close to $\pm1$; with the signs used here the data constrain $\delta K - f\,\delta\mathrm{Moho}_{00}$, which gives a correlation near $+1$, whereas the audit reported $-1$: check the parameterisation), and $K_0$ is fitted to the same elevation data the inversion uses.
\end{decision}

\section{How much each part matters}
\begin{center}
\small
\begin{tabular}{p{0.28\textwidth}p{0.3\textwidth}p{0.32\textwidth}}
\toprule
& Self-gravitating ocean & Sea-level constant\\
\midrule
Question & how sea level varies with the geoid & where sea level is relative to the columns\\
Size & \qtyrange{1}{2}{\percent} of $w$ at $\ell=2$, less above; about \qty{0.1}{\percent} of $N$ & essential: \qty{705}{m} land--ocean step without it\\
Core model & simple closed form & not needed ($b_0$)\\
Could it be deferred? & yes, stating its size; but it is already built and verified & not if the final inversion uses total topography on land\\
\bottomrule
\end{tabular}
\end{center}
If the self-gravitating ocean is neglected, two things that are not really part of it must stay: water only at sea in the plate's restoring force, and $h = w - N$, because topography is measured from sea level.

\begin{incode}{\texttt{ocean.py} and \texttt{model.py}}
\begin{itemize}[nosep]
\item \texttt{SelfGravitatingOcean}: builds $T$ from $P$, $G_w$, $C$ and $S$, applies $(I+T)^{-1}$ by the Neumann series (\texttt{\_neumann}) and its adjoint the same way; \texttt{contraction} estimates the spectral radius by power iteration.
\item \texttt{ForwardModel.sea\_level\_constant}: the least-squares $K_0$ on the fine grid with area weights.
\item \texttt{ForwardModel.uniform\_load\_operator}: $\delta K \mapsto g\,\delta K$ as the degree-0 component (the orthonormal $Y_{00}$ is $1/\sqrt{4\pi}$, so a constant $c$ has component $c\sqrt{4\pi}$).
\end{itemize}
\end{incode}
